\documentclass[11pt, a4paper]{report}

\usepackage[utf8]{inputenc}
\usepackage[T1]{fontenc}
\usepackage[french]{babel}

% Appel des packages dans le dossier preambule
\def\packagepath{../../../../../preambule}   % path du package princial

\usepackage{\packagepath/page_layout} % <--- Nouveau
\usepackage{\packagepath/config_info}
\usepackage{\packagepath/cadres_cours}
\usepackage{\packagepath/zones_reponse}
\usepackage{\packagepath/config_ds}
\usepackage{\packagepath/config_maths}

% --- PILOTAGE ---
%\def\visibleornot{visible}    % visible or invisible
\ifdefined\visibleornot\else\def\visibleornot{visible}\fi


\def\documentpath{./Sources_Latex}
\graphicspath{{./Images/}}


\def\ptsexoA{0}



\begin{document}

\def\level{Premiere}  % Classe à droite
\def\course{NSI}
\def\chaptername{La boucle WHILE} % Titre au centre
\def\docname{AP04~--~TD 2}        % Identifiant à gauche

\pagestyle{DS_LP}

\NomPrenom{}

\begin{exercice_cours}[pour chaque boucle, indiquer combien de fois elle s'exécute et les valeurs prises par \texttt{i}:]
    
\begin{minipage}{0.48\linewidth}
    \begin{CodePythontexAlt}[TaillePolice=\large, Largeur=1\linewidth,Centre,Lignes=true,PremLigne=1]{}
i = 0
while i < 5:
    print(i)
    i += 1
    \end{CodePythontexAlt}
\end{minipage}\hfill{}
\begin{minipage}{0.48\linewidth}
    \begin{CodePythontexAlt}[TaillePolice=\large, Largeur=1\linewidth,Centre,Lignes=true,PremLigne=1]{}
i = 2
while i < 10:
    print(i)
    i += 2
    \end{CodePythontexAlt}
\end{minipage}

\begin{minipage}{0.48\linewidth}
    \begin{CodePythontexAlt}[TaillePolice=\large, Largeur=1\linewidth,Centre,Lignes=true,PremLigne=1]{}
i = 0
u = 2
while u < 100:
    print(u)
    i += 1
    u = 2 ** i
    \end{CodePythontexAlt}
        
\end{minipage}\hfill{}
\begin{minipage}{0.48\linewidth}
    \begin{CodePythontexAlt}[TaillePolice=\large, Largeur=1\linewidth,Centre,Lignes=true,PremLigne=1]{}
i = 0
u = 100
while u > 1:
    print(u)
    i += 1
    u = u // 2
    \end{CodePythontexAlt}
\end{minipage}
    \end{exercice_cours}
    
    \begin{minipage}{0.48\linewidth}
        \begin{exercice_cours}[Recherche de seuil]

Une valeur initialement égale à 10 augmente de 10\% à chaque itération. Combien d'itérations sont nécessaires pour que cette valeur dépasse 100 ?
    \end{exercice_cours}

    \end{minipage}\hfill
    \begin{minipage}{0.48\linewidth}
\begin{blocvisible}[height=4.5]
    \begin{CodePythontexAlt}[TaillePolice=\large, Largeur=1\linewidth,Centre,Lignes=true,PremLigne=1]{}
def seuil():
    i = 0
    u = 10
    while u < 100:
        i += 1
        u *= 1.1
    return i
    \end{CodePythontexAlt}
\end{blocvisible}
    \end{minipage}
    
    
    \begin{minipage}{0.48\linewidth}
        \begin{exercice_cours}[Produits successifs]

Ecrire la fonction \verb|produit(seuil)| qui renvoie le plus petit entier $n$ tel que 

$1 \times 2 \times 3 \times \ldots \times n > seuil$.
    \end{exercice_cours}

    \end{minipage}\hfill
    \begin{minipage}{0.48\linewidth}
\begin{blocvisible}[height=4.5]
    \begin{CodePythontexAlt}[TaillePolice=\large, Largeur=1\linewidth,Centre,Lignes=true,PremLigne=1]{}
def produit(seuil):
    n = 1
    produit = 1
    while produit <= seuil:
        n += 1
        produit *= n
    return n
    \end{CodePythontexAlt}
\end{blocvisible}


    \end{minipage}

    \begin{minipage}{0.48\linewidth}
        \begin{exercice_cours}[Somme de puissances]

Ecrire la fonction \verb|somme(seuil)| qui renvoie le plus petit entier $n$ tel que 

$1^2 + 2^2 + 3^2 + \ldots + n^2 > seuil$.
    \end{exercice_cours}

    \end{minipage}\hfill    
    \begin{minipage}{0.48\linewidth}
\begin{blocvisible}[height=4.5] 
    \begin{CodePythontexAlt}[TaillePolice=\large, Largeur=1\linewidth,Centre,Lignes=true,PremLigne=1]{}
def somme(seuil):
    n = 1
    somme = 0
    while somme <= seuil:
        n += 1
        somme += n ** 2
    return n
    \end{CodePythontexAlt}
\end{blocvisible}
    \end{minipage}     

    
    
\medskip

\end{document}